% Nivel 2: recta, curva de 90° (radio ~1 m) y recta.
\begin{frame}[plain]
\begin{Lienzo}
\Cabecera{V · Resultados}{Nivel 2: curva de 90° de radio 1 m}
\Foto{36}{82}{290}{182}{n2_curva}
\Foto{36}{274}{290}{182}{n2_barra}
\BotonVideo{36}{470}{nivel2.mp4}{Vídeo nivel 2}
\Parrafo{196}{487}{130}{Calibri}{13}{15}{4E4B4C}{Tirada 104456, escala 1}
\Fila{346}{82}{578}{60}{PROBLEMA}{La cinta rayada de peligro junto a la curva tenía más contraste que la línea: se tapó (el PDF pide 1.5 m libres).}
\Fila{346}{150}{578}{60}{PROBLEMA}{Las dos primeras tiradas pararon 17–25 cm antes de la barra: el final a ciegas se contaba con la velocidad de mientras frenaba.}
\Fila{346}{218}{578}{60}{CAMBIO}{Velocidad medida 1.5 s después de frenar y \textbf{reajuste local de la línea hasta 0.8 m}: la parábola de toda la vista sesgaba el ángulo.}
\Figura{346}{290}{330}{170}{fig_nivel2_ajuste}
\Pie{346}{476}{330}{Error del ángulo al acercarse a la curva, frente a la IMU (3 recorridos).}
\Recuadro{692}{290}{232}{214}
\Chip{692}{290}{RESULTADO}
\CifraBarra[30]{716}{326}{200}{7 de 7}{acaban en FIN sin intervención; las 5 posteriores al arreglo, sobre la barra}
\CifraBarra[30]{716}{420}{200}{18.5–20.3 s}{a escala 1, con la parada prevista a +9…+10 cm de la barra}
\end{Lienzo}
\end{frame}
